% ==========================================================================
% Section 4: The Isomorph-Eval Framework
% 
% Compatible with neurips_2026.sty (or neurips_2025.sty / any NeurIPS template)
% Requires: amsmath, amssymb, amsthm, enumitem
%
% Preamble additions needed in main.tex:
%   \usepackage{amsmath,amssymb,amsthm}
%   \usepackage[inline]{enumitem}
%   \newtheorem{definition}{Definition}
%   \newtheorem{conjecture}{Conjecture}
%   \newtheorem{proposition}{Proposition}
% ==========================================================================

\section{The Isomorph-Eval Framework}
\label{sec:framework}

We now present the three components of Isomorph-Eval: the formal
definition of task isomorphism (\S\ref{sec:isomorphism}), the
generative engine (\S\ref{sec:engine}), and the contamination diagnostic
(\S\ref{sec:delta}).

% ------------------------------------------------------------------
\subsection{Item Decomposition and $\tau$-Isomorphism}
\label{sec:isomorphism}

\begin{definition}[Item Decomposition]
\label{def:decomposition}
An evaluation item $i$ is a tuple $(S_i, G_i, V_i)$ where:
\begin{itemize}[nosep, leftmargin=*]
\item $S_i$ is the \emph{surface representation}: the natural-language
  string presented to the model;
\item $G_i = (N_i, E_i, \tau_i, \phi_i)$ is the \emph{reasoning graph},
  a directed acyclic graph where $N_i$ is a finite set of computation
  nodes, $E_i \subseteq N_i \times N_i$ is a set of directed dependency
  edges, $\tau_i : N_i \to \mathcal{T}$ maps each node to an operation
  type from a finite taxonomy $\mathcal{T}$, and
  $\phi_i : N_i \to \Phi$ assigns a complexity weight to each node;
\item $V_i : \mathcal{A} \to \{0, 1\}$ is the \emph{verification
  function} that determines the correctness of a candidate answer.
\end{itemize}
\end{definition}

The operation taxonomy $\mathcal{T}$ is domain-specific. For
mathematical reasoning (our primary domain), we define:
\begin{equation}
\mathcal{T}_{\text{math}} = \{\textsc{assign}, \textsc{add},
\textsc{subtract}, \textsc{multiply}, \textsc{divide},
\textsc{modulo}\}.
\label{eq:taxonomy}
\end{equation}
The complexity weight $\phi(n) = (k, d, f, s)$ encodes the number of
operands $k$, the digit class $d$ (the maximum number of digits of any
operand), a Boolean $f$ indicating the presence of fractions, and a
Boolean $s$ indicating negative values. The digit class is the primary
driver of arithmetic difficulty: multiplying two 3-digit numbers is
substantially harder than multiplying two 1-digit numbers, even though
both are instances of \textsc{multiply}.

Leaf nodes are \textsc{assign} operations with no incoming edges,
representing the given quantities in the problem. Internal nodes
represent computation steps. The answer is the value at a designated
output node $n^* \in N_i$.

\begin{definition}[$\tau$-Isomorphism]
\label{def:isomorphism}
Two items $i$ and $i'$ are \emph{$\tau$-isomorphic}, written
$i \cong_\tau i'$, if and only if there exists a bijection
$\psi : N_i \to N_{i'}$ satisfying all of the following:
\begin{enumerate}[nosep, leftmargin=*, label=(\roman*)]
\item \textbf{Structural preservation:} $(n_a, n_b) \in E_i
  \iff (\psi(n_a), \psi(n_b)) \in E_{i'}$.
\item \textbf{Operation-type preservation:}
  $\forall\, n \in N_i : \tau_i(n) = \tau_{i'}(\psi(n))$.
\item \textbf{Complexity preservation:}
  $\forall\, n \in N_i : \phi_i(n) = \phi_{i'}(\psi(n))$.
\item \textbf{Semantic novelty:}
  $\mathrm{sim}(S_i, S_{i'}) < \delta$ for a threshold $\delta$, and
  $S_{i'} \notin \mathcal{C}_{\mathrm{web}}$, where
  $\mathcal{C}_{\mathrm{web}}$ denotes the set of strings present in
  publicly crawled web corpora.
\end{enumerate}
\end{definition}

\noindent
Conditions (i)--(iii) preserve the \emph{cognitive structure} of the
item: the same chain of reasoning steps, in the same dependency order,
with the same quantitative complexity at each step. Condition (iv)
ensures that the variant cannot have appeared in any training corpus,
eliminating contamination by construction. Unlike prior rephrasing
approaches~\cite{dekoninck2024constat,zhang2024gsm1k} that assess
``similar difficulty'' via heuristic comparison, $\tau$-isomorphism is a
\emph{deterministic, verifiable} structural condition: given two
reasoning graphs, one can check in polynomial time (via graph
isomorphism algorithms such as VF2~\cite{cordella2004vf2}) whether
$\psi$ exists.

\paragraph{Structural fingerprint.}
For efficient pre-screening, we define a hash
$h(G_i) = \texttt{SHA256}(|N_i|, |E_i|, \mathrm{depth}(G_i),
[\tau(n)]_{n \in \mathrm{topo}(G_i)},
[\phi(n)]_{n \in \mathrm{topo}(G_i)})$ computed over the topological
ordering of nodes. Two items can be $\tau$-isomorphic only if
$h(G_i) = h(G_{i'})$; this provides $O(1)$ rejection of non-isomorphic
pairs before invoking the full VF2 check.


% ------------------------------------------------------------------
\subsection{The Reasoning Structure Hypothesis}
\label{sec:hypothesis}

The utility of $\tau$-isomorphism for contamination detection rests on
the following conjecture linking graph structure to IRT parameters.

\begin{conjecture}[Reasoning Structure Hypothesis]
\label{conj:rsh}
Let item $i$ have IRT parameters $(a_i, b_i)$ estimated from a
sufficiently large calibration panel of $J$ systems. If
$i \cong_\tau i'$, then:
\begin{equation}
|b_i - b_{i'}| \leq \epsilon_b
\quad \text{and} \quad
|a_i - a_{i'}| \leq \epsilon_a
\label{eq:rsh}
\end{equation}
for tolerance bounds $\epsilon_b, \epsilon_a$ that depend on the
granularity of the operation taxonomy $\mathcal{T}$ and the complexity
weight space $\Phi$.
\end{conjecture}

\noindent
The intuition is that difficulty $b$ in the 2PL model reflects the
ability level at which 50\% of systems succeed. If two items require
the same chain of reasoning operations with the same structural
dependencies and the same quantitative load at each node, then a system
that cannot perform the required computation chain will fail both, and a
system that can will succeed at both---regardless of whether the
entities are named ``Janet'' or ``Tom\'{a}s'' and whether the units are
``eggs'' or ``honey jars.''

This is a conjecture, not an axiom: surface features \emph{can} affect
difficulty through secondary channels (unfamiliar vocabulary, numeric
magnitude within digit class, culturally specific framing). The
complexity weight $\phi$ mitigates the dominant source of such leakage
(digit class), and the empirical calibration protocol
(Section~\ref{sec:delta}) provides a falsifiability mechanism.


% ------------------------------------------------------------------
\subsection{The Isomorphic Engine}
\label{sec:engine}

The Isomorphic Engine is a four-stage pipeline that automatically
generates verified $\tau$-isomorphic variants. Given an original item
$(S, G, V)$, it produces a variant $(S', G', V')$ such that
$G \cong_\tau G'$ and $V'$ is automatically determined by forward
execution of $G'$.

\paragraph{Stage 1: Parser ($S \to G, V$).}
We employ a hybrid strategy. A rule-based annotation parser extracts the
computation DAG from GSM8K's \texttt{<<expr=result>>} annotations,
achieving 57.2\% coverage on the test set (755/1{,}319 problems) with
zero LLM involvement. For the remaining items, we use LLM structured
extraction via Pydantic-enforced schemas and the Anthropic tool-use
API, with the LLM's output validated by forward graph
execution.\footnote{If the graph's computed answer does not match the
ground-truth answer, the parse is rejected and retried---we never trust
the LLM's extraction blindly.}

\paragraph{Stage 2: Mutator ($G \to G'$).}
The mutator operates exclusively on the \emph{free dimensions} of the
item: entity bindings (e.g., ``Janet'' $\to$ ``Tom\'{a}s''), numeric
leaf values (preserving digit class), units (``dollars'' $\to$
``euros''), and context theme. The \emph{locked invariants}---node set,
edge set, operation types $\tau$, and complexity weights $\phi$---are
never modified. After mutation, the answer $V'$ is recomputed by forward
execution of $G'$, requiring no human annotation.

\paragraph{Stage 3: Surface Generator ($G' \to S'$).}
We use a \emph{template-first} architecture. A deterministic template
containing all mathematical content is generated from $G'$, then
optionally polished by an LLM. The LLM's system prompt strictly forbids
changing any number or logical relationship; its role is exclusively
that of a surface realizer, not a reasoner. The prompt enforces six
explicit constraints, and the output is post-checked for number
preservation.

\paragraph{Stage 4: Verifier ($S' \to G'' \cong_\tau G'$).}
The critical safety mechanism is \emph{round-trip verification}: the
generated surface text $S'$ is re-parsed into a reasoning graph $G''$
using the same parser as Stage~1, and $G''$ is checked against $G'$ for
$\tau$-isomorphism. This catches any corruption introduced by the
surface generator, including subtle errors such as implicit assumptions,
reordered operations, or added conditions. Items that fail round-trip
verification are discarded. The verifier also runs seven static checks:
topology match, operation-type sequence match, complexity-weight
sequence match, structural fingerprint match, number presence, answer
consistency, and digit-class preservation.


% ------------------------------------------------------------------
\subsection{The Contamination Delta}
\label{sec:delta}

\begin{definition}[Item-Level Contamination Delta]
\label{def:item_delta}
For model $M$ evaluated on original item $i$ with $T$ trials and on $K$
$\tau$-isomorphic variants $\{i'_1, \ldots, i'_K\}$ with $T$ trials
each:
\begin{equation}
\hat{p}_i^{\,\mathrm{orig}} = \frac{1}{T} \sum_{t=1}^{T} X_{i,t},
\qquad
\hat{p}_i^{\,\mathrm{iso}} = \frac{1}{KT} \sum_{k=1}^{K}
\sum_{t=1}^{T} X_{i'_k, t},
\qquad
\delta_i(M) = \hat{p}_i^{\,\mathrm{orig}} -
\hat{p}_i^{\,\mathrm{iso}}
\label{eq:item_delta}
\end{equation}
where $X_{i,t} \in \{0,1\}$ is the binary outcome of model $M$ on item
$i$ at trial $t$ (with temperature $> 0$ to induce sampling variance).
Under $H_0$ (no contamination), $\mathbb{E}[\delta_i] = 0$ because
$i \cong_\tau i'_k$ implies identical IRT parameters (Conjecture~\ref{conj:rsh}).
\end{definition}

\begin{definition}[IRT-Weighted Contamination Delta]
\label{def:model_delta}
For model $M$ evaluated on benchmark $B = \{i_1, \ldots, i_N\}$:
\begin{equation}
\Delta_{\mathrm{contam}}^{\mathrm{IRT}}(M, B) = \frac{1}{N}
\sum_{i=1}^{N} \hat{a}_i \cdot \delta_i(M)
\label{eq:delta_irt}
\end{equation}
where $\hat{a}_i$ is the discrimination parameter of item $i$ estimated
from the clean (isomorphic) data via 2PL IRT. Items with higher
discrimination receive greater weight, because a performance gap on a
highly discriminating item---one that effectively separates high-ability
from low-ability systems---constitutes stronger evidence of memorization
than a gap on a non-discriminating item.
\end{definition}

\paragraph{Hypothesis test via Differential Item Functioning.}
We ground our statistical test in the psychometric framework of
Differential Item Functioning (DIF)~\cite{lord1980,holland1993},
which detects items that function differently across groups matched on
latent ability. In our setting, the two ``groups'' are evaluation
conditions: Condition A (model sees original items) and Condition B
(model sees $\tau$-isomorphic variants).

We fit separate 2PL IRT models to each condition:
\begin{align}
\text{Model}_{\mathrm{orig}}&: \quad
P(X_{ji} = 1) = \sigma\!\left(a_i^{\mathrm{orig}}
(\theta_j^{\mathrm{orig}} - b_i^{\mathrm{orig}})\right),
\label{eq:irt_orig}
\\
\text{Model}_{\mathrm{iso}}&: \quad
P(X_{ji'} = 1) = \sigma\!\left(a_i^{\mathrm{iso}}
(\theta_j^{\mathrm{iso}} - b_i^{\mathrm{iso}})\right).
\label{eq:irt_iso}
\end{align}
Under $H_0$ (no contamination):
$b_i^{\mathrm{orig}} = b_i^{\mathrm{iso}}$ and
$a_i^{\mathrm{orig}} = a_i^{\mathrm{iso}}$ for all $i$.
Under $H_1$ (contamination): $b_i^{\mathrm{orig}} <
b_i^{\mathrm{iso}}$ for contaminated items (the original appears
easier because the model has memorized it).

\paragraph{Per-item Wald test.}
For each item $i$, we test $H_0 : b_i^{\mathrm{orig}} =
b_i^{\mathrm{iso}}$ using:
\begin{equation}
z_i = \frac{\hat{b}_i^{\,\mathrm{orig}} -
\hat{b}_i^{\,\mathrm{iso}}}
{\mathrm{SE}(\hat{b}_i^{\,\mathrm{orig}} -
\hat{b}_i^{\,\mathrm{iso}})}
\label{eq:wald}
\end{equation}
with Bonferroni correction: item $i$ is flagged as ``memorized'' if
$|z_i| > z_{\alpha / 2N}$. The global test uses a likelihood ratio
statistic $\Lambda = -2[\ell_0 - \ell_1] \sim \chi^2(N)$, comparing
the constrained model ($b^{\mathrm{orig}} = b^{\mathrm{iso}}$
$\forall\, i$) against the unconstrained model.


% ------------------------------------------------------------------
\subsection{Ability Decomposition and the Three-Body Problem}
\label{sec:threebody}

The IRT-weighted contamination delta enables a formal decomposition of
observed ability:

\begin{proposition}[Ability Decomposition]
\label{prop:decomp}
Under the 2PL model, the observed ability $\theta^{\mathrm{obs}}_j$ of
model $j$ on the original benchmark can be decomposed as:
\begin{equation}
\theta^{\mathrm{obs}}_j = \theta^{\mathrm{true}}_j +
\Delta\theta^{\mathrm{contam}}_j
\label{eq:decomp}
\end{equation}
where $\theta^{\mathrm{true}}_j$ is the ability estimated from
isomorphic items (clean), and
$\Delta\theta^{\mathrm{contam}}_j = \theta^{\mathrm{obs}}_j -
\theta^{\mathrm{true}}_j$ is the ability inflation due to memorization.
\end{proposition}

\noindent
This decomposition is valid because $\tau$-isomorphic items share the
same $(a, b)$ parameters (by Conjecture~\ref{conj:rsh}); the difference
$\theta^{\mathrm{obs}} - \theta^{\mathrm{true}}$ is therefore
attributable to item-specific memorization, not to any legitimate
ability difference.

\paragraph{The Three-Body Problem.}
In our prior work~\cite{kang2026efsl}, we established that ranking
error under simple averaging is well described by:
\begin{equation}
1 - \rho_{\mathrm{avg}} = \gamma_0 + \gamma_1 S_c + \gamma_2 D_c +
\gamma_3 (S_c \times D_c) + \varepsilon
\label{eq:efsl}
\end{equation}
where $S$ is sparsity, $D$ is the difficulty gap, and centered variables
$S_c = S - \bar{S}$, $D_c = D - \bar{D}$. We now extend this to
include contamination as a third axis:
\begin{equation}
1 - \rho = f(S, D, C)
\label{eq:threebody}
\end{equation}
where $C$ measures the mean differential contamination across models.
The compound failure is worst when all three co-occur: a contaminated
model tested on easy items in a sparse evaluation matrix receives
\emph{triple-inflated} ability estimates. IRT alone corrects for $S$
and $D$ (via joint estimation of $\theta_j$ and $b_i$), but cannot
correct for $C$ because it cannot distinguish genuine ability from
memorization on the same items. Isomorphic items alone correct for $C$
(by providing clean test items) but do not correct for $S$ and $D$ if
the resulting evaluation matrix is still sparse and heterogeneous.
\textbf{Only the combination}---IRT estimation on isomorphic data---corrects
all three axes simultaneously.

\paragraph{Diagnostic archetypes.}
Based on the pattern of ($\Delta_{\mathrm{contam}}$, $\Delta\theta$,
DIF), we classify models into four archetypes
(Table~\ref{tab:archetypes}):
\emph{Robust Reasoner}
($\Delta_{\mathrm{contam}} \leq 0.02$, no flagged items);
\emph{Pure Memorizer}
($\Delta_{\mathrm{contam}} \geq 0.10$, $>$30\% items flagged);
\emph{Syntactic Matcher}
($0.03 \leq \Delta_{\mathrm{contam}} < 0.10$, mixed DIF);
and \emph{Latent Contaminator}
(low per-item $\Delta$ but cross-benchmark mismatch, capturing the
Goodhart-in-checkpoint-selection phenomenon identified by
\citet{zhang2024gsm1k}).
